\section{Introduction}
\label{sec:intro}

Software often has to choose among fixed options: route a support ticket, classify an intent, or check a claim against evidence.
A chat model can do this, but its answer must be generated and parsed.
Jev-style decision models \citep{jev} instead return a probability distribution over the options, so a new task needs only a prompt and a list of options.
LLM-as-Jev \citep{li2026llmasjev} shows that an unmodified decoder language model works as such a model: it reads numbered options and answers with a number.
Fine-tuning it with KL anchors helps most for small models and for tasks with many options.

A decision needs almost no generation, while decoders are built for it.
Multi-query and grouped-query attention \citep{mqa, gqa} shrink the KV cache for decoding, and the parameters they save can go into larger MLP layers.
\citet{prefillonly} show that these layers cause memory spikes in workloads that only read their input.
An encoder keeps no KV cache, which is one reason to try it; at our scale, \cref{sec:cost} finds no peak-memory advantage.

An encoder also reads the input in both directions.
A classifier head would fix the label set, but a masked language model can score any option by whether it fills a mask next to the option with ``yes'' or ``no''.
Recent encoders such as ModernBERT \citep{modernbert} and UltraBERT \citep{ultrabert} read the long inputs that decisions over documents need.

Reading in both directions also helps with evidence.
A decoder serving many questions about one document usually reads the document first, so that its KV cache can be reused, as with prefix caching in vLLM \citep{vllm} and RadixAttention in SGLang \citep{sglang}.
The cached document states then do not depend on the question.
Our encoder reads the question and the document together.

Our contributions are:
\begin{itemize}\setlength{\itemsep}{2pt}\setlength{\topsep}{2pt}
  \item \textbf{An encoder decision model.}
    UltraBERT-380M with \emph{one mask per option}, scoring \emph{every option in one forward pass} with \emph{no new head}, for any number of options.
    It is fine-tuned on the LLM-as-Jev training mixture with Nimble added (\cref{sec:data}).
  \item \textbf{Training recipes.}
    \emph{Four ways to score and normalise} the mask predictions, one following the DeepSeek-V3 expert gate \citep{deepseekv3}.
    \emph{Scoring by ``yes'' minus ``no''} gives the larger gain; \emph{sigmoid normalisation} adds \DiffUbSigUbYesNoAbs{} test points, mostly on MuSiQue.
  \item \textbf{A comparison with decoders under matched data.}
    Against Qwen3-0.6B \citep{qwen3} fully fine-tuned on the same examples, \emph{every encoder recipe is more accurate} (\TestUbList{}\% to \TestUbSig{}\% against \TestQwFull{}\% on a held-out test split) and \emph{faster} (\MsUb{} against \MsQwFull{}\,ms per decision), but \emph{not lighter} in peak memory.
    Two reproduced LoRA models \citep{lora} serve as references.
  \item \textbf{Predicting supporting evidence.}
    An \emph{auxiliary evidence head} in the same forward pass, trained on human-marked labels and labels selected with a reward adapted from SelfCite \citep{selfcite}.
    On the softmax encoder it \emph{ranks annotated sentences higher than ContextCite and AT2} \citep{contextcite, at2}.
\end{itemize}
